\section{Generalized linear models. Lasso, Ridge. Generalized additive models. }
\subsection{Generalized linear models}


	Lasso and Ridge from the book.\\
	Lasso formula, Ridge formula, explanation \\
	Въпрос: Как Ridge и Lasso fit-ват в общите линейни модели или
	общите адитивни модели?
	
	
	
	Да припомним няколко регресии, които сме виждали:\\
	\begin{itemize}
		\item линейна - $\mathbb{E}(Y|X_1, \dots, X_p) = \beta_0 + \beta_1 X_1 + \dots + \beta_p X_p $
		\item логистична -  $ \mathbb{E}(Y|X_1, \dots, X_p) = \frac{e^{\beta_0 + \beta_1 X_1 + \dots + \beta_p X_p}}{1+ e^{\beta_0 + \beta_1 X_1 + \dots + \beta_p X_p}} $
		\item Поансонова -  $ \mathbb{E}(Y|X_1, \dots, X_p) = e^{\beta_0 + \beta_1 X_1 + \dots + \beta_p X_p} $
 	\end{itemize}
	
	Общият линеен модел изглежда по следния начин(където $\eta$ е $"$свързваща фунцкия$"$):\\
	$$\eta(\mathbb{E}(Y|X_1, \dots, X_p)) = \beta_0 + \beta_1 X_1 + \dots + \beta_p X_p $$
	
	lab 6 за ridge lasso \\
	Change $\lambda$ for Lasso example.
	
	
	
	Подробно разписване на оптимизационната задача. Идея за множество от бъдещите задачи
	се разглеждат като оптимизационни.
	
\newpage	
